\chapter{Penerapan Teori Grup}\label{ch:b3:group-theory-applied}

Karbon dioksida hanya menyusun beberapa ratus molekul dari setiap sejuta
molekul udara, sedangkan nitrogen dan oksigen hampir seluruh sisanya; namun
karbon dioksidalah, bukan nitrogen atau oksigen, yang menyerap cahaya
inframerah yang dipancarkan tanah yang hangat ke antariksa. Sebuah molekul
hanya menyerap cahaya inframerah melalui vibrasi yang mengubah momen
dipolnya, dan apakah suatu vibrasi berbuat demikian ditentukan oleh simetri
semata: satu-satunya vibrasi nitrogen tidak dapat melakukannya, sedangkan
dua dari empat vibrasi karbon dioksida dapat. Bab ini mengubah tabel-tabel
\omterm{def:b3:point-groups:representation}{karakter} \cref{ch:b3:point-groups} menjadi alat: bab ini mereduksi
representasi, menyusun orbital-orbital teradaptasi simetri, menghitung dan
menandai vibrasi, serta menurunkan aturan seleksi spektroskopi inframerah dan
Raman.

\begin{recall}
\Cref{ch:b3:point-groups} mendefinisikan \omterm{def:b3:point-groups:operation}{operasi simetri}, \omterm{def:b3:point-groups:point-group}{grup titik}, kelas,
representasi, \omterm{def:b3:point-groups:representation}{karakter}, dan \omterm{def:b3:point-groups:irrep}{tabel karakter}. Jilid Tahun 2 menyusun orbital
molekul \ce{H2O}, \ce{NH3} dan \ce{CH4} dari orbital-orbital fragmen dan
kombinasi teradaptasi simetri orbital $1s$ hidrogen, dengan menamainya
menggunakan label-label yang diturunkan di sini. Jilid Tahun 1 membaca
spektrum inframerah dalam bilangan gelombang dan mendefinisikan
polarisabilitas sebuah molekul.
\end{recall}

\section{Mereduksi sebuah representasi}

Setiap himpunan fungsi atau vektor yang melekat pada sebuah molekul membawa
sebuah representasi \omterm{def:b3:point-groups:point-group}{grup titiknya}. Sebagian besar di antaranya tereduksi.

\begin{definition}[Representasi tereduksi, representasi simetris total]\label{def:b3:group-theory-applied:reducible}
Yang disebut \emph{representasi tereduksi}\index{representasi tereduksi}
adalah representasi yang semua matriksnya dapat dibawa, dengan satu
perubahan basis, ke bentuk diagonal blok yang sama; representasi itu lalu
merupakan jumlah \omterm{def:b3:point-groups:irrep}{representasi tak tereduksi}, ditulis
$\Gamma = \sum_in_i\Gamma_i$. \omterm{def:b3:point-groups:irrep}{Representasi tak tereduksi} yang semua
\omterm{def:b3:point-groups:representation}{karakternya} 1 disebut \emph{representasi simetris
total}\index{representasi simetris total} ($A_1$, $A_{1g}$, $A_1'$\dots).
\end{definition}

\begin{theorem}[Teorema ortogonalitas besar]\label{thm:b3:group-theory-applied:great-orthogonality}
Untuk dua \omterm{def:b3:point-groups:irrep}{representasi tak tereduksi} $\Gamma_i$, $\Gamma_j$ berdimensi $d_i$,
$d_j$ dari sebuah grup berorde $h$, yang ditulis dengan matriks uniter,
\[
\sum_R D_i(R)_{mn}^*\,D_j(R)_{m'n'} = \frac{h}{d_i}\,\delta_{ij}\,\delta_{mm'}\,\delta_{nn'}.
\]
\end{theorem}
\admitted

Buktinya, dengan lema-lema Schur, dibahas dalam kuliah yang lebih lanjut;
yang dipakai kimia adalah akibat-akibatnya bagi \omterm{def:b3:point-groups:representation}{karakter}, dan setiap \omterm{def:b3:point-groups:irrep}{tabel
karakter} dalam buku ini diperiksa terhadapnya.

\begin{corollary}[Keortogonalan karakter]\label{cor:b3:group-theory-applied:characters}
Dengan $g_c$ banyaknya operasi dalam kelas $c$,
\[
\sum_cg_c\,\chi_i(c)^*\chi_j(c) = h\,\delta_{ij},
\]
dan banyaknya \omterm{def:b3:point-groups:irrep}{representasi tak tereduksi} sama dengan banyaknya kelas, dengan
$\sum_id_i^2 = h$.
\end{corollary}

\begin{proof}
Ambil $m = n$ dan $m' = n'$ dalam teorema itu lalu jumlahkan atas $m$ dan
$m'$: ruas kiri menjadi $\sum_R\chi_i(R)^*\chi_j(R)$, ruas kanan $\frac{h}{d_i}\delta_{ij}
\sum_{m}\delta_{mm} = h\delta_{ij}$; kelompokkan operasi-operasinya menurut kelas.
Dengan demikian, baris-baris tabel adalah vektor-vektor ortogonal dalam ruang
yang dimensinya sama dengan banyaknya kelas, sehingga \omterm{def:b3:point-groups:irrep}{representasi tak
tereduksi} paling banyak sebanyak kelasnya; kesamaan dan $\sum d_i^2 = h$
diperoleh dengan menerapkan teorema itu pada representasi reguler
(latihan 10).
\end{proof}

\begin{theorem}[Rumus reduksi]\label{thm:b3:group-theory-applied:reduction}
Sebuah representasi $\Gamma$ dengan \omterm{def:b3:point-groups:representation}{karakter} $\chi(c)$ memuat \omterm{def:b3:point-groups:irrep}{representasi tak
tereduksi} $\Gamma_i$ tepat
\[
n_i = \frac1h\sum_cg_c\,\chi(c)\,\chi_i(c)^*
\]
kali. Inilah \emph{rumus reduksi}\index{rumus reduksi}.
\end{theorem}

\begin{proof}
Karena $\Gamma = \sum_jn_j\Gamma_j$, \omterm{def:b3:point-groups:representation}{karakternya} adalah $\chi(c) = \sum_jn_j\chi_j(c)$.
Kalikan dengan $g_c\chi_i(c)^*$, jumlahkan atas kelas-kelas, dan gunakan
akibat di atas: $\sum_cg_c\chi(c)\chi_i(c)^* = \sum_jn_jh\delta_{ij} = hn_i$.
\end{proof}

\begin{example}[Orbital-orbital hidrogen pada amonia]\label{ex:b3:group-theory-applied:nh3}
Ambil ketiga orbital $1s$ hidrogen \ce{NH3} sebagai basis. $E$ membiarkan
ketiganya di tempat ($\chi = 3$), $C_3$ memindahkan ketiganya ($\chi = 0$),
sebuah $\sigma_v$ membiarkan satu di tempat dan menukar dua lainnya
($\chi = 1$). Dengan tabel $C_{3v}$:
$n_{A_1} = \frac16(3 + 0 + 3\times1) = 1$, $n_{A_2} = \frac16(3 + 0 - 3) = 0$,
$n_E = \frac16(6 + 0 + 0) = 1$. Jadi $\Gamma_{\mathrm H} = A_1 + E$: ketiga
hidrogen memberikan satu kombinasi simetris total dan sepasang kombinasi
terdegenerasi.
\end{example}

\begin{method}[Mereduksi sebuah representasi]\label{met:b3:group-theory-applied:reduce}
\begin{enumerate}
  \item Untuk setiap kelas, carilah \omterm{def:b3:point-groups:representation}{karakternya}: untuk himpunan fungsi yang
        berpusat pada atom, banyaknya fungsi yang tetap di tempat, masing-masing
        dihitung dengan tanda yang diperolehnya.
  \item Terapkan \omterm{thm:b3:group-theory-applied:reduction}{rumus reduksi} dengan baris setiap \omterm{def:b3:point-groups:irrep}{representasi tak tereduksi},
        dengan memberi bobot pada setiap kelas sebesar ukurannya.
  \item Periksalah: hasilnya bilangan bulat tak negatif, dan $\sum_in_id_i$
        sama dengan \omterm{def:b3:point-groups:representation}{karakter} di bawah $E$ (ukuran basisnya).
\end{enumerate}
\end{method}

\section{Orbital teradaptasi simetri}

\begin{definition}[Operator proyeksi]\label{def:b3:group-theory-applied:projection}
Yang disebut \emph{operator proyeksi}\index{operator proyeksi} ke
\omterm{def:b3:point-groups:irrep}{representasi tak tereduksi} $\Gamma_i$ adalah
\[
\hat P_i = \frac{d_i}{h}\sum_R\chi_i(R)^*\,\hat R,
\]
dengan $\hat R$ mengenakan operasi $R$ pada sebuah fungsi.
\end{definition}

\begin{proposition}[Kerja operator proyeksi]\label{prop:b3:group-theory-applied:projection}
Bila dikenakan pada sembarang fungsi $f$, $\hat P_if$ bernilai nol atau
merupakan fungsi yang bertransformasi seperti $\Gamma_i$ (sebuah kombinasi
teradaptasi simetri); untuk $\Gamma_i$ berdimensi satu, fungsi itu tidak
berubah, kecuali tanda $\chi_i(S)$, oleh setiap operasi $S$.
\end{proposition}

\begin{proof}[Bukti parsial]
Untuk $d_i = 1$: $\hat S\hat P_if = \frac1h\sum_R\chi_i(R)\hat S\hat Rf$. Ambil $R' = SR$,
yang menjelajahi seluruh grup ketika $R$ menjelajahinya; $\chi_i(R) = \chi_i(S^{-1}R') =
\chi_i(S)\chi_i(R')$ untuk representasi berdimensi satu (\omterm{def:b3:point-groups:representation}{karakternya} adalah
matriks itu sendiri). Jadi $\hat S\hat P_if = \chi_i(S)\hat P_if$. Kasus umumnya
memakai teorema ortogonalitas besar dan diterima tanpa bukti.
\end{proof}

\begin{method}[Menyusun kombinasi teradaptasi simetri]\label{met:b3:group-theory-applied:salc}
\begin{enumerate}
  \item Pilihlah satu fungsi dari himpunan itu, $h_1$, dan tabulasikan fungsi
        yang menjadi bayangannya oleh setiap operasi.
  \item Untuk setiap \omterm{def:b3:point-groups:irrep}{representasi tak tereduksi}, kalikan setiap bayangan
        dengan \omterm{def:b3:point-groups:representation}{karakter} operasinya, jumlahkan, lalu normalisasikan (dengan
        mengabaikan tumpang-tindih antarfungsi atom).
  \item Untuk representasi terdegenerasi, proyeksikan fungsi kedua untuk
        memperoleh anggota kedua, lalu ortogonalkan.
\end{enumerate}
\end{method}

\begin{example}[Kombinasi-kombinasi pada amonia dan air]\label{ex:b3:group-theory-applied:salcs}
Pada \ce{NH3}, $\hat P_{A_1}h_1 \propto h_1 + h_2 + h_3$ (semua \omterm{def:b3:point-groups:representation}{karakter} 1). Untuk $E$
(\omterm{def:b3:point-groups:representation}{karakter} $2, -1, 0$), $\hat P_Eh_1 \propto 2h_1 - h_2 - h_3$; memproyeksikan $h_2$
lalu mengortogonalkannya memberikan pasangannya, $h_2 - h_3$. Setelah
dinormalisasi: $a_1 =
\frac{1}{\sqrt3}(h_1 + h_2 + h_3)$, $e = \frac{1}{\sqrt6}(2h_1 - h_2 - h_3)$,
$\frac{1}{\sqrt2}(h_2 - h_3)$. Pada \ce{H2O}, yang diletakkan pada bidang $xz$,
$h_1 + h_2$ adalah $a_1$ dan $h_1 - h_2$ adalah $b_1$ (kombinasi itu berganti
tanda oleh $C_2$ dan oleh $\sigma_v'(yz)$, yang mempertukarkan
hidrogen-hidrogennya). Jika molekul diletakkan pada bidang $yz$, seperti yang
disukai sebagian buku, label $b_1$ dan $b_2$ bertukar di seluruh uraian.
\end{example}

\begin{omfigure}
\begin{tikzpicture}[font=\small]
  \foreach \x/\lab/\ca/\cb/\cc in {0/{$a_1$}/1/1/1, 4/{$e$}/2/-1/-1, 8/{$e$}/0/1/-1}{
    \begin{scope}[xshift=\x cm]
      \foreach \a/\c in {90/\ca,210/\cb,330/\cc}{
        \pgfmathsetmacro{\r}{0.18+0.13*abs(\c)}
        \ifdim \c pt > 0pt \fill[omDef!55, draw=omDef] (\a:1.0) circle (\r);
        \else \ifdim \c pt < 0pt \fill[omThm!25, draw=omThm] (\a:1.0) circle (\r);
        \else \draw[gray, densely dotted] (\a:1.0) circle (0.12); \fi\fi
      }
      \node[aN, minimum size=4.6mm] at (0,0) {};
      \node[anchor=north] at (0,-1.45) {\lab};
    \end{scope}}
  \node[anchor=north, font=\footnotesize] at (0,-1.85) {$h_1 + h_2 + h_3$};
  \node[anchor=north, font=\footnotesize] at (4,-1.85) {$2h_1 - h_2 - h_3$};
  \node[anchor=north, font=\footnotesize] at (8,-1.85) {$h_2 - h_3$};
\end{tikzpicture}

{\small Kombinasi teradaptasi simetri ketiga orbital $1s$ hidrogen pada
amonia, dilihat sepanjang sumbu $C_3$ ($h_1$ di atas). Biru: koefisien positif,
merah: negatif, dengan ukuran setiap cakram tumbuh bersama koefisiennya;
titik-titik: nol.}
\end{omfigure}

\begin{example}[Enam ligan di sekitar sebuah logam]\label{ex:b3:group-theory-applied:ml6}
Untuk enam orbital $\sigma$ ligan yang mengarah ke sebuah logam dari
titik-titik sudut oktahedron, \omterm{def:b3:point-groups:representation}{karakter-karakternya} pada kelas-kelas $O_h$
($E$, $8C_3$, $6C_2$, $6C_4$, $3C_2$, $i$, $6S_4$, $8S_6$, $3\sigma_h$, $6\sigma_d$)
adalah $6, 0, 0, 2, 2, 0, 0, 0, 4,
2$, dan reduksinya memberikan $\Gamma_\sigma = A_{1g} + E_g + T_{1u}$. Orbital $s$
logam ($a_{1g}$), $d_{z^2}$ dan $d_{x^2-y^2}$ ($e_g$) serta $p$ ($t_{1u}$)
menemukan pasangan ligan; orbital $d_{xy}$, $d_{xz}$, $d_{yz}$ ($t_{2g}$) tidak
menemukannya dan tetap nonikatan dalam $\sigma$ --- asal-usul pemisahan
$t_{2g}$/$e_g$ dalam jilid Tahun 2, yang dikembangkan dalam
\cref{ch:b3:organometallic-bonding,ch:b3:complex-spectra-magnetism}.
\end{example}

\begin{omfigure}
\begin{tikzpicture}[font=\small, yscale=0.95]
  % O orbitals (left), MOs (centre), H combinations (right); schematic energies
  \foreach \y/\l in {0/{$2s\ (a_1)$}}{\draw[omstate] (0,\y) -- (1.2,\y); \node[anchor=east] at (-0.05,\y) {\l};}
  \draw[omstate] (0,3.0) -- (0.36,3.0) (0.42,3.0) -- (0.78,3.0) (0.84,3.0) -- (1.2,3.0);
  \node[anchor=east] at (-0.05,3.0) {$2p\ (b_1, a_1, b_2)$};
  \draw[omstate] (6.8,3.9) -- (8.0,3.9); \node[anchor=west] at (8.05,3.9) {$h_1 - h_2\ (b_1)$};
  \draw[omstate] (6.8,3.6) -- (8.0,3.6); \node[anchor=west] at (8.05,3.5) {$h_1 + h_2\ (a_1)$};
  \foreach \y/\l in {-0.6/{$2a_1$}, 1.6/{$1b_1$}, 2.4/{$3a_1$}, 3.0/{$1b_2$}, 5.6/{$4a_1$}, 6.2/{$2b_1$}}{
    \draw[omstate] (3.4,\y) -- (4.6,\y); \node[anchor=west] at (4.65,\y) {\l};}
  \foreach \y in {-0.6,1.6,2.4,3.0}{\node[font=\scriptsize] at (4.0,\y+0.16) {$\uparrow\downarrow$};}
  \draw[densely dotted, gray] (1.2,0) -- (3.4,-0.6) (1.2,3.0) -- (3.4,1.6) (1.2,3.0) -- (3.4,2.4) (1.2,3.0) -- (3.4,3.0)
     (6.8,3.6) -- (4.6,-0.6) (6.8,3.9) -- (4.6,1.6) (6.8,3.6) -- (4.6,2.4) (1.2,3.0) -- (3.4,6.2) (6.8,3.9) -- (4.6,6.2) (6.8,3.6) -- (4.6,5.6);
  \node[anchor=north] at (0.6,-1.0) {O};
  \node[anchor=north] at (4.0,-1.0) {\ce{H2O}};
  \node[anchor=north] at (7.4,-1.0) {2 H};
  \draw[->] (-2.3,-0.9) -- (-2.3,6.4) node[above] {$E$};
\end{tikzpicture}

{\small Orbital molekul valensi air (energi skematis, molekul pada bidang $xz$).
Hanya orbital-orbital bersimetri sama yang bercampur: kombinasi hidrogen
$b_1$ dengan $2p_x$, kombinasi $a_1$ dengan $2s$ dan $2p_z$; $2p_y$ ($b_2$)
tidak mempunyai pasangan dan tetap menjadi pasangan elektron bebas
nonikatan, orbital terisi tertinggi.}
\end{omfigure}

\section{Modus vibrasi}

\begin{definition}[Modus normal]\label{def:b3:group-theory-applied:normal-mode}
Yang disebut \emph{modus normal}\index{modus normal} sebuah molekul adalah
vibrasi kolektif yang semua atomnya berosilasi dengan frekuensi yang sama dan
sefase, sepanjang vektor eigen Hessian berbobot massa
(\cref{prop:b3:computational-chemistry:hessian}). Setiap modus normal
bertransformasi seperti sebuah \omterm{def:b3:point-groups:irrep}{representasi tak tereduksi} \omterm{def:b3:point-groups:point-group}{grup titiknya}.
\end{definition}

\begin{proposition}[Representasi semua gerakan]\label{prop:b3:group-theory-applied:gamma-3n}
Pasangkan tiga vektor perpindahan $x$, $y$, $z$ pada setiap atom. Oleh sebuah
operasi $R$, \omterm{def:b3:point-groups:representation}{karakter} representasi berdimensi $3N$ ini adalah
\[
\chi_{3N}(R) = (\text{banyaknya atom yang tetap di tempat}) \times \chi_{xyz}(R),
\]
dengan $\chi_{xyz} = 1 + 2\cos\theta$ untuk rotasi sebesar $\theta$ dan
$-1 + 2\cos\theta$ untuk rotasi tak sejati (jadi 3 untuk $E$, $-1$ untuk $C_2$,
0 untuk $C_3$, 1 untuk $\sigma$, $-3$ untuk $i$).
\end{proposition}

\begin{proof}
Atom yang dipindahkan ke posisi lain membawa vektor-vektornya keluar dari
diagonal matriks: atom itu tidak menyumbang apa pun pada jejak. Atom yang
tetap di tempat ketiga vektornya ditransformasikan oleh matriks $3\times3$
dari $R$, yang jejaknya dihitung dalam kerangka dengan $z$ sepanjang sumbu:
$\cos\theta + \cos\theta + 1$ untuk rotasi, dengan entri terakhir menjadi $-1$
untuk rotasi tak sejati.
\end{proof}

\begin{proposition}[Menghitung vibrasi]\label{prop:b3:group-theory-applied:mode-count}
$\Gamma_{\mathrm{vib}} = \Gamma_{3N} - \Gamma_{\mathrm{trans}} - \Gamma_{\mathrm{rot}}$, dengan
$\Gamma_{\mathrm{trans}}$ representasi $(x, y, z)$ dan $\Gamma_{\mathrm{rot}}$ representasi
$(R_x, R_y, R_z)$; molekul nonlinear mempunyai $3N - 6$ vibrasi, molekul
linear $3N - 5$.
\end{proposition}

\begin{proof}
Ke-$3N$ perpindahan itu menggambarkan setiap gerakan: tiga translasi pusat
massa, tiga rotasi (dua untuk molekul linear, karena berputar terhadap
sumbunya sendiri tidak memindahkan inti mana pun), dan vibrasi-vibrasinya.
Subruang-subruang ini tidak bercampur oleh operasi-operasi itu, sehingga
representasinya dapat dijumlahkan.
\end{proof}

\begin{example}[Air]\label{ex:b3:group-theory-applied:water}
% ledger: vib:H2O.sym, vib:H2O.bend, vib:H2O.asym
\ce{H2O} dalam $C_{2v}$, molekul pada bidang $xz$. Atom yang tetap di tempat: 3,
1, 3, 1; $\chi_{xyz}$: $3, -1, 1, 1$; jadi $\chi_{3N} = 9, -1, 3, 1$. Reduksinya
memberikan $3A_1 +
A_2 + 3B_1 + 2B_2$. Setelah translasi $A_1 + B_1 + B_2$ dan rotasi $A_2 + B_1 + B_2$
dikurangkan: $\Gamma_{\mathrm{vib}} = 2A_1 + B_1$. Kedua modus $a_1$ adalah
vibrasi ulur simetris (\qty{3657}{cm^{-1}}) dan vibrasi tekuk
(\qty{1595}{cm^{-1}}); modus $b_1$ adalah vibrasi ulur antisimetris
(\qty{3756}{cm^{-1}}).
\end{example}

\begin{omfigure}
\begin{tikzpicture}[font=\small, >=Stealth]
  \foreach \x/\name/\lab in {0/sym/{$\nu_1$, $a_1$, ulur simetris}, 4.3/bend/{$\nu_2$, $a_1$, tekuk}, 8.6/asym/{$\nu_3$, $b_1$, ulur antisimetris}}{
    \begin{scope}[xshift=\x cm]
      \node[aO] (o\name) at (0,0.35) {};
      \node[aH] (a\name) at (-0.8,-0.25) {};
      \node[aH] (b\name) at (0.8,-0.25) {};
      \begin{scope}[on background layer]\draw[bond] (o\name.center) -- (a\name.center) (o\name.center) -- (b\name.center);\end{scope}
      \node[anchor=north, align=center, font=\footnotesize] at (0,-0.9) {\lab};
    \end{scope}}
  % symmetric stretch: H outward along the bonds, O up a little
  \draw[->, omThm, thick] (-0.95,-0.36) -- (-1.42,-0.71);
  \draw[->, omThm, thick] (0.95,-0.36) -- (1.42,-0.71);
  \draw[->, omThm, thick] (0,0.74) -- (0,0.98);
  % bend: each H moves perpendicular to its bond, closing the angle; O up
  \draw[->, omThm, thick] (4.3-0.92,-0.36) -- (4.3-0.62,-0.76);
  \draw[->, omThm, thick] (4.3+0.92,-0.36) -- (4.3+0.62,-0.76);
  \draw[->, omThm, thick] (4.3,0.74) -- (4.3,0.98);
  % antisymmetric stretch: one H out, the other in, O sideways
  \draw[->, omThm, thick] (8.6-0.95,-0.36) -- (8.6-1.42,-0.71);
  \draw[->, omThm, thick] (8.6+1.02,0.0) -- (8.6+0.70,0.24);
  \draw[->, omThm, thick] (8.6-0.05,0.80) -- (8.6+0.25,0.80);
\end{tikzpicture}

{\small Ketiga \omterm{def:b3:group-theory-applied:normal-mode}{modus normal} air (panah: perpindahan, tidak berskala). Kedua
modus $a_1$ mempertahankan simetri penuh molekul; modus $b_1$ berganti tanda
oleh $C_2$.}
\end{omfigure}

\begin{example}[Empat molekul lagi]\label{ex:b3:group-theory-applied:table}
Prosedur yang sama (dihitung dan diperiksa dalam data gambar bab ini)
memberikan:
\begin{center}\small
\begin{tabular}{llll}
\hline
molekul & grup & $\Gamma_{\mathrm{vib}}$ & modus \\ \hline
\ce{NH3} & $C_{3v}$ & $2A_1 + 2E$ & 6 \\
\ce{CH4} & $T_d$ & $A_1 + E + 2T_2$ & 9 \\
\ce{CO2} & $D_{\infty h}$ (via $D_{2h}$) & $\Sigma_g^+ + \Sigma_u^+ + \Pi_u$ ($A_g + B_{1u} + B_{2u} + B_{3u}$) & 4 \\
\ce{BF3} & $D_{3h}$ & $A_1' + 2E' + A_2''$ & 6 \\
\ce{XeF4} & $D_{4h}$ & $A_{1g} + B_{1g} + B_{2g} + A_{2u} + B_{2u} + 2E_u$ & 9 \\
\ce{SF6} & $O_h$ & $A_{1g} + E_g + T_{2g} + 2T_{1u} + T_{2u}$ & 15 \\ \hline
\end{tabular}
\end{center}
Untuk molekul linear, grup tak berhingganya diganti dengan \omterm{def:b3:point-groups:class}{subgrupnya}
$D_{2h}$, tempat vibrasi tekuk terdegenerasi $\Pi_u$ terpecah menjadi
$B_{2u} + B_{3u}$.
\end{example}

\section{Aturan seleksi}

\begin{definition}[Hasil kali langsung]\label{def:b3:group-theory-applied:direct-product}
Yang disebut \emph{hasil kali langsung}\index{hasil kali langsung}
$\Gamma_i\otimes\Gamma_j$ dua representasi adalah representasi yang dibawa oleh
hasil kali fungsi-fungsi basis keduanya; \omterm{def:b3:point-groups:representation}{karakternya} adalah hasil kali
$\chi_i(R)\chi_j(R)$.
\end{definition}

\begin{theorem}[Integral yang lenyap]\label{thm:b3:group-theory-applied:vanishing-integral}
Integral $\int f_1f_2f_3\,\dd\tau$ atas seluruh ruang hanya dapat bernilai tak
nol jika \omterm{def:b3:group-theory-applied:direct-product}{hasil kali langsung} $\Gamma_1\otimes\Gamma_2\otimes\Gamma_3$ memuat
\omterm{def:b3:group-theory-applied:reducible}{representasi simetris total}.
\end{theorem}

\begin{proof}
Sebuah \omterm{def:b3:point-groups:operation}{operasi simetri} hanya menandai ulang titik-titik ruang, sehingga nilai
integral tidak berubah bila integrannya ditransformasikan oleh $R$ mana pun.
Rata-ratakan integran yang ditransformasikan itu atas grupnya: rata-rata
komponen-komponen integran yang termasuk dalam setiap \omterm{def:b3:point-groups:irrep}{representasi tak
tereduksi} adalah $\frac1h\sum_R\hat R(\cdot)$, yaitu proyeksi ke \omterm{def:b3:group-theory-applied:reducible}{representasi
simetris total} (\cref{def:b3:group-theory-applied:projection} dengan semua
$\chi = 1$). Setiap komponen lain berata-rata nol; jika hasil kali itu tidak
memuat bagian simetris total, integralnya nol.
\end{proof}

\begin{definition}[Aktif IR, aktif Raman]\label{def:b3:group-theory-applied:activity}
Sebuah vibrasi fundamental disebut \emph{aktif IR}\index{aktif IR} jika dapat
menyerap cahaya inframerah, dan \emph{aktif Raman}\index{aktif Raman} jika
muncul dalam spektrum Raman (\cref{ch:b3:rovibrational-spectroscopy}).
\end{definition}

\begin{corollary}[Aturan seleksi inframerah dan Raman]\label{cor:b3:group-theory-applied:ir-raman}
Sebuah fundamental (dari tingkat dasar, yang simetris total, ke satu kuantum
sebuah modus bersimetri $\Gamma$) hanya \omterm{def:b3:group-theory-applied:activity}{aktif IR} jika $\Gamma$ adalah
representasi $x$, $y$ atau $z$, dan hanya \omterm{def:b3:group-theory-applied:activity}{aktif Raman} jika $\Gamma$ adalah
representasi sebuah fungsi kuadratik ($x^2$, $xy$, \dots).
\end{corollary}

\begin{proof}
Intensitas serapan melibatkan $\int\psi_0\,\mu_k\,\psi_1\dd\tau$ dengan
komponen-komponen dipol $\mu_k$, yang bertransformasi seperti $x$, $y$, $z$;
$\psi_0$ simetris total dan $\psi_1$ bertransformasi seperti $\Gamma$. Menurut
teorema itu, integralnya hanya dapat tak nol jika $\Gamma\otimes\Gamma_{\mu_k}$
memuat $A_1$, dan hal ini hanya terjadi jika $\Gamma =
\Gamma_{\mu_k}$ (hasil kali dua \omterm{def:b3:point-groups:irrep}{representasi tak tereduksi} memuat \omterm{def:b3:group-theory-applied:reducible}{representasi
simetris total} hanya jika keduanya sama, untuk \omterm{def:b3:point-groups:representation}{karakter} real). Hamburan Raman
melibatkan komponen-komponen polarisabilitas $\alpha_{kl}$, yang
bertransformasi seperti hasil kali $kl$.
\end{proof}

\begin{proposition}[Aturan eksklusi timbal balik]\label{prop:b3:group-theory-applied:mutual-exclusion}
Dalam molekul yang mempunyai \omterm{def:b3:point-groups:inversion}{pusat inversi}, tidak ada fundamental yang
sekaligus \omterm{def:b3:group-theory-applied:activity}{aktif IR} dan \omterm{def:b3:group-theory-applied:activity}{aktif Raman}: inilah \emph{aturan eksklusi timbal
balik}\index{aturan eksklusi timbal balik}.
\end{proposition}

\begin{proof}
Dengan adanya \omterm{def:b3:point-groups:inversion}{pusat inversi}, setiap \omterm{def:b3:point-groups:irrep}{representasi tak tereduksi} bersifat $g$
atau $u$. Koordinat $x$, $y$, $z$ berganti tanda oleh $i$ (ketiganya $u$);
fungsi-fungsi kuadratik tidak ($g$). Sebuah modus hanya \omterm{def:b3:group-theory-applied:activity}{aktif IR} jika $u$, dan
hanya \omterm{def:b3:group-theory-applied:activity}{aktif Raman} jika $g$.
\end{proof}

\begin{example}[Karbon dioksida dan efek rumah kaca]\label{ex:b3:group-theory-applied:co2}
% ledger: vib:CO2.sym, vib:CO2.bend, vib:CO2.asym, ir:CO2.asym.int, ir:CO2.bend.int
\ce{CO2} mempunyai \omterm{def:b3:point-groups:inversion}{pusat inversi}. Vibrasi ulur simetrisnya ($\Sigma_g^+$,
\qty{1333}{cm^{-1}}) hanya \omterm{def:b3:group-theory-applied:activity}{aktif Raman}; vibrasi ulur antisimetris ($\Sigma_u^+$,
\qty{2349}{cm^{-1}}) dan vibrasi tekuk ($\Pi_u$, \qty{667}{cm^{-1}}) hanya \omterm{def:b3:group-theory-applied:activity}{aktif
IR}. Vibrasi tekuk menyerap di sekitar \qty{15}{\micro m}, dekat puncak emisi
termal Bumi: itulah sebabnya karbon dioksida merupakan gas rumah kaca.
\ce{N2} dan \ce{O2} hanya mempunyai satu vibrasi, $\Sigma_g^+$: tidak \omterm{def:b3:group-theory-applied:activity}{aktif IR}.
\end{example}

\begin{omfigure}
\begin{tikzpicture}
\begin{axis}[omir, width=0.92\linewidth, height=4.4cm, xmin=400, xmax=2600,
  ymin=0, ymax=1.3, ytick=\empty, xlabel={bilangan gelombang / cm$^{-1}$},
  ylabel={sinyal}, legend cell align=left,
  legend style={font=\small, at={(0.98,0.95)}, anchor=north east, draw=none}]
\addplot[omDef, thick] table[x=nu, y=ir] {figdata/out/bachelor-3-co2-spectra.dat};
\addlegendentry{serapan inframerah}
\addplot[omThm, thick, dashed] table[x=nu, y=raman] {figdata/out/bachelor-3-co2-spectra.dat};
\addlegendentry{hamburan Raman}
\node[font=\small, anchor=south] at (axis cs:2349,1.02) {$\Sigma_u^+$};
\node[font=\small, anchor=south] at (axis cs:667,0.1) {$\Pi_u$};
\node[font=\small, anchor=south, omThm] at (axis cs:1333,1.02) {$\Sigma_g^+$};
\end{axis}
\end{tikzpicture}

{\small Spektrum sintetis gas \ce{CO2} (posisi pita dan intensitas relatif
inframerah dari nilai-nilai terukur; bentuknya skematis, tanpa struktur
rotasi). Eksklusi timbal balik: tidak ada pita yang muncul pada keduanya.}
\end{omfigure}

\begin{method}[Meramalkan spektrum vibrasi]\label{met:b3:group-theory-applied:vibrations}
\begin{enumerate}
  \item Carilah \omterm{def:b3:point-groups:point-group}{grup titiknya}; hitunglah $\chi_{3N}$ dari atom-atom yang
        tetap di tempat.
  \item Reduksikan, lalu kurangkan translasi dan rotasi: $\Gamma_{\mathrm{vib}}$.
  \item Dari kolom-kolom kanan tabel, tandai setiap modus sebagai \omterm{def:b3:group-theory-applied:activity}{aktif IR}
        ($x$, $y$, $z$) dan/atau \omterm{def:b3:group-theory-applied:activity}{aktif Raman} (kuadratik); modus yang tidak
        termasuk keduanya disebut senyap.
  \item Hitunglah pita-pitanya: satu untuk setiap modus aktif, dengan modus
        terdegenerasi memberikan satu pita.
\end{enumerate}
\end{method}

\begin{method}[Menghitung vibrasi ulur CO]\label{met:b3:group-theory-applied:co-count}
Pada karbonil logam, ambil $n$ vektor ikatan C--O sebagai basis (atom yang
tetap di tempat dihitung 1, selainnya tidak dihitung): reduksikan untuk
memperoleh $\Gamma_{\mathrm{CO}}$, lalu terapkan aturan IR dan Raman. Banyaknya
pita kuat di sekitar \qtyrange{1850}{2150}{cm^{-1}} mengidentifikasi isomernya.
\end{method}

\begin{example}[Cis dan trans]\label{ex:b3:group-theory-applied:cis-trans}
cis-\ce{ML4(CO)2}, $C_{2v}$: \omterm{def:b3:point-groups:representation}{karakter} $2, 0, 2, 0$, $\Gamma_{\mathrm{CO}} = A_1 + B_1$,
dua pita IR. trans-\ce{ML4(CO)2}, $D_{4h}$: $\Gamma_{\mathrm{CO}} = A_{1g} + A_{2u}$,
satu pita IR ($A_{2u}$) dan satu pita Raman ($A_{1g}$). Satu pita CO saja dalam
spektrum inframerah berarti trans.
\end{example}

\section{Penerapan pada transisi elektronik}

Teorema yang sama berlaku untuk transisi elektronik: transisi elektronik
antara keadaan $\Psi_1$ dan $\Psi_2$ hanya diizinkan jika $\Gamma_1\otimes\Gamma_2$
memuat representasi $x$, $y$ atau $z$, dan cahayanya terpolarisasi sepanjang
sumbu itu.

\begin{example}[Transisi-transisi air dan formaldehida]\label{ex:b3:group-theory-applied:electronic}
Pada air ($C_{2v}$, bidang $xz$), mempromosikan sebuah elektron dari $1b_2$
(pasangan elektron bebas) ke $4a_1$ memberikan keadaan tereksitasi bersimetri
$B_2\otimes A_1 = B_2$, yaitu $y$: diizinkan, terpolarisasi tegak lurus pada
bidang molekul. Pada formaldehida ($C_{2v}$), transisi $n\to\pi^*$ berlangsung
dari pasangan elektron bebas oksigen yang sebidang dengan molekul ($b_1$) ke
orbital $\pi^*$, yang tersusun dari orbital-orbital $p$ yang tegak lurus pada
bidang itu ($b_2$; molekul pada bidang $xz$, C=O sepanjang $z$):
$B_1\otimes B_2 = A_2$, yang bukan $x$, $y$, maupun $z$: terlarang oleh
simetri, itulah sebabnya pitanya (\cref{ch:b3:electronic-spectroscopy}) begitu
lemah.
\end{example}

\begin{history}[Wigner, Bethe, dan simetri molekul]
Teori grup masuk ke mekanika kuantum melalui karya Eugene Wigner tentang
spektrum atom (1927--31) dan makalah Hans Bethe tentang pemecahan suku-suku
atom dalam kristal (1929). E. Bright Wilson menerapkannya pada vibrasi
molekul pada tahun 1934; Robert Mulliken memberikan label-label yang masih
dipakai untuk orbital dan keadaan.
\end{history}

\begin{inthelab}[Inframerah dan Raman berdampingan]
Spektrometer inframerah mengukur cahaya yang diserap sampel; spektrometer
Raman menyinari sampel dengan laser dan menganalisis cahaya hamburan lemah
pada panjang gelombang yang bergeser. Merekam keduanya pada senyawa yang sama
adalah uji klasik untuk pusat simetri: jika tidak ada pita yang berimpit,
molekul itu kemungkinan sentrosimetris. Laser Raman termasuk kelas 3B atau 4:
berkasnya tertutup dan para operatornya memakai kacamata pelindung yang
sesuai dengan panjang gelombangnya.
\end{inthelab}

\section{Latihan}

\begin{exercise}[$\star$]\label{exo:b3:group-theory-applied:1}
Reduksikan representasi $C_{2v}$ dengan \omterm{def:b3:point-groups:representation}{karakter} $(5, -1, 3, 1)$ pada $(E, C_2,
\sigma_v(xz), \sigma_v'(yz))$. Periksalah dimensinya.
\end{exercise}

\begin{exercise}[$\star$]\label{exo:b3:group-theory-applied:2}
\ce{SO2} berbentuk bengkok, seperti air. Berikan $\Gamma_{\mathrm{vib}}$ dan
tentukan modus mana yang \omterm{def:b3:group-theory-applied:activity}{aktif IR} dan \omterm{def:b3:group-theory-applied:activity}{aktif Raman}.
\end{exercise}

\begin{exercise}[$\star$]\label{exo:b3:group-theory-applied:3}
Dalam $C_{2v}$, hitunglah \omterm{def:b3:group-theory-applied:direct-product}{hasil kali langsung} $B_1\otimes B_2$ dan
$A_2\otimes B_1$, serta $E\otimes E$ dalam $C_{3v}$ (reduksikan yang terakhir).
\end{exercise}

\begin{exercise}[$\star$]\label{exo:b3:group-theory-applied:4}
% ledger: vib:CH4.nu1, vib:CH4.nu2, vib:CH4.nu3, vib:CH4.nu4
\ce{CH4} mempunyai fundamental pada 2917 ($a_1$), 1534 ($e$), 3019 dan 1306
(keduanya $t_2$) \unit{cm^{-1}}. Manakah yang muncul dalam spektrum
inframerah, dan manakah dalam spektrum Raman?
\end{exercise}

\begin{exercise}[$\star\star$]\label{exo:b3:group-theory-applied:5}
Untuk \ce{BF3}, turunkan $\Gamma_{\mathrm{vib}} = A_1' + 2E' + A_2''$ dari atom-atom
yang tetap di tempat, lalu hitunglah pita IR dan pita Raman. Modus mana yang
hanya terlihat dalam salah satu spektrum, dan modus mana yang terlihat dalam
keduanya?
\end{exercise}

\begin{exercise}[$\star\star$]\label{exo:b3:group-theory-applied:6}
Berapa banyak pita IR dan pita Raman yang diperlihatkan \ce{SF6}? Modus mana
yang senyap? Mengapa tidak ada pita yang muncul dalam kedua spektrum?
\end{exercise}

\begin{exercise}[$\star\star$]\label{exo:b3:group-theory-applied:7}
Kenakan \omterm{def:b3:group-theory-applied:projection}{operator proyeksi} pada $h_2$ untuk representasi $E$ dari $C_{3v}$ dan
tunjukkan bahwa mengortogonalkan hasilnya terhadap $2h_1 - h_2 - h_3$
memberikan $h_2 - h_3$.
\end{exercise}

\begin{exercise}[$\star\star$]\label{exo:b3:group-theory-applied:8}
Ramalkan banyaknya pita CO dalam IR untuk fac- dan mer-\ce{ML3(CO)3}, serta
untuk \ce{M(CO)6}.
\end{exercise}

\begin{exercise}[$\star\star$]\label{exo:b3:group-theory-applied:9}
Tunjukkan bahwa keenam orbital $\sigma$ ligan sebuah kompleks oktahedral
merentang $A_{1g} + E_g + T_{1u}$, dan sebutkan orbital-orbital logam mana yang
dapat berkombinasi dengan masing-masing.
\end{exercise}

\begin{exercise}[$\star\star\star$]\label{exo:b3:group-theory-applied:10}
Representasi reguler mempunyai $\chi(E) = h$ dan $\chi(R) = 0$ untuk $R \ne E$.
Tunjukkan dengan \omterm{thm:b3:group-theory-applied:reduction}{rumus reduksi} bahwa representasi itu memuat setiap
\omterm{def:b3:point-groups:irrep}{representasi tak tereduksi} $\Gamma_i$ tepat $d_i$ kali, lalu simpulkan
$\sum_id_i^2 = h$.
\end{exercise}

\begin{exercise}[$\star\star\star$]\label{exo:b3:group-theory-applied:11}
Asetilena \ce{HC#CH} linear. Dengan bekerja dalam $D_{2h}$, carilah
$\Gamma_{\mathrm{vib}}$, kelompokkan kembali ke dalam label-label $D_{\infty h}$,
dan ramalkan spektrum IR dan Raman molekul itu.
\end{exercise}

\begin{exercise}[$\star\star\star$]\label{exo:b3:group-theory-applied:12}
Pada air (bidang $xz$), manakah di antara promosi satu elektron berikut yang
memberikan transisi yang diizinkan, dan dengan polarisasi apa:
$1b_2 \to 4a_1$, $3a_1 \to 4a_1$, $1b_2 \to 2b_1$, $1b_1 \to 4a_1$?
\end{exercise}

\section{Soal: Cis atau Trans? Membaca Spektrum Karbonil}

\begin{problem}[{Soal akhir pekan --- grup titik dan representasi vibrasi
ulur CO empat isomer, pita inframerah dan Raman masing-masing, serta identifikasi
sebuah produk dari spektrum terukurnya}]
\label{pb:b3:group-theory-applied:1}
% ledger: ct:C2v, ct:C3v, ct:D4h
Sebuah logam M oktahedral membawa ligan-ligan karbonil dan ligan-ligan lain
L (setiap L dihitung sebagai satu titik). Empat isomer ditinjau: cis- dan
trans-\ce{ML4(CO)2}, serta fac- dan mer-\ce{ML3(CO)3}. Anggaplah ligan-ligan L
tidak menurunkan simetri melebihi apa yang ditentukan oleh posisinya. Sebuah
sintesis \ce{ML3(CO)3} memberikan produk yang spektrum inframerahnya
memperlihatkan tiga pita kuat pada 2048, 1965 dan \qty{1938}{cm^{-1}} (data
soal); spektrum Raman produk itu memperlihatkan tiga pita pada posisi yang hampir
sama.

\medskip\noindent\textbf{Bagian I --- \omterm{def:b3:point-groups:point-group}{Grup titik}.}
\begin{enumerate}[beginpenalty=10000]
  \item Gambarlah keempat isomer itu pada sebuah oktahedron.
  \item Berikan \omterm{def:b3:point-groups:point-group}{grup titik} cis-\ce{ML4(CO)2}.
  \item \omterm{def:b3:point-groups:point-group}{Grup titik} trans-\ce{ML4(CO)2}.
  \item \omterm{def:b3:point-groups:point-group}{Grup titik} fac-\ce{ML3(CO)3}.
  \item \omterm{def:b3:point-groups:point-group}{Grup titik} mer-\ce{ML3(CO)3}.
  \item Manakah di antara keempatnya yang mempunyai \omterm{def:b3:point-groups:inversion}{pusat inversi}?
\end{enumerate}

\medskip\noindent\textbf{Bagian II --- Vibrasi ulur CO.}
\begin{enumerate}[resume, beginpenalty=10000]
  \item Jelaskan mengapa vektor-vektor ikatan C--O membawa sebuah representasi
        yang di dalamnya sebuah operasi menyumbang 1 untuk setiap CO yang tetap
        di tempat.
  \item Carilah $\Gamma_{\mathrm{CO}}$ isomer cis.
  \item Untuk isomer trans.
  \item Untuk isomer fac.
  \item Untuk isomer mer.
  \item Periksalah setiap dimensi terhadap banyaknya ligan CO.
\end{enumerate}

\medskip\noindent\textbf{Bagian III --- Keaktifan.}
\begin{enumerate}[resume, beginpenalty=10000]
  \item Hitunglah vibrasi ulur CO yang \omterm{def:b3:group-theory-applied:activity}{aktif IR} untuk setiap isomer.
  \item Hitunglah yang \omterm{def:b3:group-theory-applied:activity}{aktif Raman}.
  \item Isomer mana yang memperlihatkan \omterm{prop:b3:group-theory-applied:mutual-exclusion}{aturan eksklusi timbal balik}?
  \item Gambarkan modus \omterm{def:b3:group-theory-applied:activity}{aktif IR} isomer trans: apakah kedua CO mengulur sefase?
  \item Pada isomer fac, mengapa ketiga CO hanya memberikan dua pita?
  \item Apa yang akan diberitahukan oleh hitungan pita sebanyak 1 (IR)?
\end{enumerate}

\medskip\noindent\textbf{Bagian IV --- Produknya.}
\begin{enumerate}[resume, beginpenalty=10000]
  \item Isomer mana yang ditunjukkan oleh spektrum inframerah produk itu?
  \item Apakah spektrum Raman itu konsisten?
  \item Pita tertinggi adalah vibrasi ulur sefase gugus-gugus CO. Termasuk
        \omterm{def:b3:point-groups:irrep}{representasi tak tereduksi} mana pita itu pada isomer tersebut?
  \item Mungkinkah itu campuran kedua isomer? Pengukuran lanjutan apa yang memutuskannya?
  \item Nyatakan hasilnya: banyaknya vibrasi ulur CO yang \omterm{def:b3:group-theory-applied:activity}{aktif IR} pada isomer
        mer, dibandingkan dengan isomer fac.
\end{enumerate}
\end{problem}
